\section{Methods}

\subsection{Study design, estimands and analysis status}

This study was a secondary computational analysis of two redistributed dense
Vina-family score matrices (Docking-44 and DOCKSTRING-58), two redistributed
experimental kinase panels (DAVIS and PKIS2), and the checksum-verified official
Anastassiadis HotSpot workbook, which is fetched but not redistributed. No new
docking or experimental measurements were performed. The primary matrix-level estimand
was the Pearson target-correlation map on an explicitly stated finite ligand
support, before and after removal of additive ligand and target effects.  Map
agreement was Spearman correlation between corresponding strict-upper-triangle
edges.  The experimental boundary analysis used the same map estimand on a
fixed panel of 21 kinases and, for the DAVIS/PKIS2 same-support comparison, on
exact matched ligands. Panel-transfer utility was signed experimental
co-response coverage on those same fixed target labels; assay values were not
used to select Vina panels.

The analyses were not preregistered and no prospective power calculation was
performed.  Molecular weight was selected after preliminary inspection; its
quartile comparison, threshold sweep and continuum analysis are exploratory.
The descriptor-specificity, pilot-recovery, panel-compression,
chemical-group-holdout, experimental-reliability and three-panel transfer
analyses were post hoc revision analyses. Randomization probabilities and repeated-sample ranges are
therefore used to characterize a specified finite-support procedure, not as
confirmatory error control over the analysis history.  Ligand libraries, target
panels and assay systems are treated as fixed observed supports; no interval in
this study is a confidence interval for a superpopulation of targets.  The
cross-panel QAP is conditional on the 21 named target labels.

\subsection{Public score matrices, preprocessing and overlap audit}

Docking-44 was read from the frozen project export
\path{data/frozen/df_final_v4.csv.gz}, derived from Nikitin et al.
\cite{nikitin2025}.  It contains 12,651 ligand rows and 44 Vina-GPU~2.0 target
columns \cite{ding2023}.  Numeric score cells were clipped above at zero and
the 8,159 missing cells (1.466\% of 556,644 cells, distributed across 3,354
rows) were replaced by the mean of their target column for the primary
full-support analysis.  The frozen matrix already contained no positive value;
111 cells were exactly zero, so the clipping operation did not alter it.

DOCKSTRING-58 was read from the DOCKSTRING v1 table
\texttt{data/frozen/}\allowbreak\texttt{dockstring-}\allowbreak\texttt{dataset.tsv.gz}
\cite{dockstring2022}.  The source
contains 260,155 molecules and 58 AutoDock Vina target columns
\cite{vina2010}.  The primary support retained the 260,060 rows complete across
all 58 targets, excluding 95 rows containing any of 339 missing cells.  Of the
retained scores, 5,393 were positive and were clipped to zero.  Retaining the
positive values and retaining all 260,155 rows under target-mean or
target-median imputation were recorded preprocessing sensitivities in the
public producer.

The two matrices were generated by different source projects, but were not
treated as statistically independent samples.  Chemical overlap was audited at
three levels: full Standard InChIKey, its first 14 characters (the connectivity
block), and non-empty canonical Bemis--Murcko scaffold \cite{murcko1996}.  For
Docking-44, Standard InChIKeys were recomputed from the analysis SMILES.  For
DOCKSTRING, they were read from the frozen, row-aligned identity contract
\path{data/frozen/dockstring_identity_contract_2026-08-03.csv.gz}; this contract
parses the raw source SMILES with the pinned RDKit version and records the raw
Standard InChIKey and connectivity block without applying a fragment-parent or
uncharging transform.  Empty Murcko scaffolds were not counted in the
cross-matrix scaffold intersection.  Target identities and the curated
same-receptor-structure indicator were taken from the audited mapping table.
These audits produced the reported counts of shared
identities, connectivity blocks, scaffolds, nine mapped targets and three
shared receptor structures; they did not remove any row from the full spectral
analyses.

The analysis starts at these published, frozen score matrices.  It reruns all
matrix transformations and comparisons reported here, but does not rerun
receptor preparation, protonation, search-box definition, pose generation or
the upstream docking campaigns.  Receptor structures, configurations and
upstream commands for DOCKSTRING remain part of its source release.  Thus the
study audits dependence in released scores and does not reproduce poses,
absolute affinities or docking accuracy.

\subsection{Target maps, centring and spectral summaries}

For an $N\times P$ finite score matrix $X$, the raw map was the sample Pearson
correlation matrix of its $P$ target columns.  The residual surface was the
two-way-centred matrix
$\Gamma_{ij}=X_{ij}-\bar X_{i\cdot}-\bar X_{\cdot j}+\bar X_{\cdot\cdot}$, and
the residual map was the target correlation matrix of $\Gamma$.  Correlation
places target columns on a common unit-variance scale using sample standard
deviations (denominator $N-1$).  Because correlation already removes column
means, the raw-to-residual map change is operationally the effect of removing
each ligand's row mean.  Two-way centring imposes zero row sums and hence one
zero spectral mode.

The eigenspectrum was calculated from a symmetrized target correlation matrix;
for both $X$ and $\Gamma$, correlation construction standardizes every target
to unit sample variance after the stated transformation, so the primary
raw--residual PR contrast is normalization matched.
numerically negative eigenvalues were clipped to zero.  Participation-ratio
(PR) dimension was
\[
 r_{\mathrm{PR}}=\frac{(\sum_{k=1}^{P}\lambda_k)^2}
 {\sum_{k=1}^{P}\lambda_k^2}.
\]
For a $P$-target correlation matrix, $\sum_k\lambda_k=P$ and
$\sum_k\lambda_k^2=P+2\sum_{j<k}r_{jk}^2$, giving the exact identity
\[
r_{\mathrm{PR}}=\frac{P}{1+(P-1)\overline{r_{jk}^{2}}},
\]
where the bar averages unique off-diagonal target pairs.  PR is therefore a
monotone re-expression of mean squared off-diagonal target correlation, not an
algebraic rank or an independent test statistic.  Entropy effective rank, PC1
variance fraction and the number of
components reaching 90\% cumulative variance were retained as secondary
spectral summaries \cite{mazzucato2016,roy2007}.  For the shared-axis diagnostic,
$X$ was column-standardized with sample standard deviations, the leading
eigenvector was oriented to have positive loading sum, and its cosine with the
unit-normalized all-ones vector was calculated.

The covariance sensitivity replaced correlation by the sample covariance
matrix of $X$ or $\Gamma$ (denominator $N-1$) and applied the same eigenvalue and
PR formula.  It intentionally retained target-specific variances; no target
standardization was applied in this sensitivity.  Agreement between any two
maps was Spearman correlation between their corresponding $P(P-1)/2$ strict
upper-triangle entries, with average ranks for ties.  Sign agreement,
root-mean-square edge error and
deterministic top-edge overlaps were reported where specified below.

\subsection{Transformation-matched residual nulls}

Each matrix used one deterministic simple-random support of 12,000 ligands,
sampled without replacement (NumPy seeds 60 for Docking-44 and 61 for
DOCKSTRING-58).  The sorted zero-based row indices were encoded as little-endian
64-bit integers and SHA-256 hashed.  All four null families within a dataset
used that identical fingerprinted support and 500 replicates.

In the additive Gaussian null, the sampled surface was decomposed into its
grand mean, ligand mean minus the grand mean, and target mean minus the grand
mean.  Independent normal noise was generated for every cell, with each target
column's standard deviation set to that target's empirical two-way-residual
sample standard deviation.  Raw and re-centred residual PR were then
recalculated.  In the empirical residual-column null, each target column of the
observed two-way residual was independently permuted across ligands, after
which the simulated matrix was two-way centred again and its PR calculated.
This preserves each target's residual marginal at the permutation step, but not
cross-target correspondence or ligand-specific residual magnitude.  In the
row-norm null, independent standard-normal $P$-vectors were generated for each
ligand, their within-row means were removed, and each vector was rescaled to
the observed Euclidean norm of that ligand's residual row.  This preserves the
support, target count, zero row sums and every row norm exactly, but not target
marginals, target variances or alignment.

The molecular-weight-conditioned null operated on processed raw scores rather
than on residuals.  RDKit molecular weights were divided by empirical quantiles
into 1, 5, 10 or 20 bins; duplicate quantile edges were removed and exact ties
were assigned together with \texttt{searchsorted(side=left)}.  Within every bin
and target, raw scores were independently permuted across ligands.  Row means
were then removed and target correlations recomputed.  At the permutation step
this preserves the exact target-specific empirical score distribution within
each coarse size bin, including non-Gaussian shape, but destroys cross-target
ligand correspondence and does not preserve continuous within-bin trends or
residual row norms.  For 5, 10 and 20 bins, a size-matched random partition
used the same bin counts after randomly assigning bin labels to ligands.  Ten
bins were primary; the other counts tested granularity.  Besides residual PR,
the null recorded PC1 fraction, mean absolute and root-mean-square correlation,
maximum absolute edge, Spearman agreement with the observed map, sign agreement
and top-edge overlap.

Null summaries are Monte Carlo distributions under these four specified data-
generating mechanisms.  Central 90\% and 95\% simulation ranges were calculated
by empirical quantiles.  The lower-tail Monte Carlo probability was
$(1+\#\{r_{\mathrm{PR}}^{(b)}\leq r_{\mathrm{PR}}^{\mathrm{obs}}\})/(500+1)$.
The ranges and probabilities do not test docking accuracy, biological
correctness or transport beyond the fixed panels.

\subsection{Molecular-weight domains and chemical-support controls}

RDKit \texttt{Descriptors.MolWt} was calculated after parsing the frozen
analysis SMILES \cite{rdkit2025}.  Docking-44 used all 12,651 primary processed
rows.  DOCKSTRING used a score-blind simple-random sample of 15,000 complete
rows without replacement (primary seed 71); the complete-row index set is
fingerprinted in the result bundle.  Five additional fixed supports (seeds 72,
73, 74, 75 and 20260809) assessed DOCKSTRING support sensitivity.

For the quartile analysis, empirical 0.25 and 0.75 quantiles were calculated and
all boundary ties were retained: low rows satisfied
$\mathrm{MW}\leq q_{0.25}$ and high rows satisfied
$\mathrm{MW}\geq q_{0.75}$.  Raw and row-centred maps were estimated separately
in the two bands.  Two controls distinguished a domain change from ordinary
finite-support or group-disjoint variation.  Equal-size row-random halves were
non-overlapping; if a band contained an odd number of rows, one row from that
permutation was unused.  For the chemical-group control, complete groups were
assigned
within ten stable rank strata of group-median molecular weight.  Within each
stratum the group order was permuted, and groups were assigned greedily to the
currently smaller ligand-count half; the initial side for exact ties was random
and then alternated.  No group crossed halves.  The same row-random and
MW-stratified whole-group splits were repeated separately inside the low and
high bands.

Docking-44 grouping used the frozen source column
\texttt{Butina\_clusters}; the strict public-core producer did not regenerate
those labels, so their fingerprint recipe is not assumed from the label name.
DOCKSTRING grouping used canonical RDKit Bemis--Murcko scaffold SMILES.  Every
acyclic DOCKSTRING molecule was deliberately assigned its own source-row label
\texttt{ACYCLIC\_SINGLETON:<source-row>}, including acyclic molecules with the
same connectivity, so that it could not leak across group-disjoint halves.
This singleton rule is specific to the DOCKSTRING chemical-domain controls.

There were 200 repetitions per control.  The base seed was 20260812 for
Docking-44 and 20261812 for DOCKSTRING.  Within each dataset, low- and high-band
row-random controls used base seed plus 100 and 101, and low- and high-band
group-disjoint controls used base seed plus 200 and 201.  Central 95\% repeated-
split ranges describe composition sensitivity, not sampling confidence
intervals.

Threshold robustness used exact equal-size lower and upper rank tails at
fractions 0.10, 0.15, 0.20, 0.25, 0.30, 0.35 and 0.40.  Rows were sorted by
molecular weight and then source-row index, making boundary ties deterministic;
each tail contained $\lfloor fN\rfloor$ rows.  For each fraction, 100 equal-$N$
random-disjoint controls used seeds 20310810 (Docking-44) and 20311810
(DOCKSTRING).  The continuum analysis divided the same supports into ten
non-overlapping equal-$N$ rank windows; the highest-ranked $N\bmod 10$ remainder
rows were excluded.  Spearman correlation between the 45 pairwise differences
in median molecular weight and $1-$map agreement was descriptive because the
45 comparisons reuse ten maps.  Molecular weight is correlated with other
chemical properties, and none of these analyses identifies a causal molecular-
weight effect.

\subsection{Descriptor-axis specificity}

As a post hoc specificity analysis, the same low-versus-high support contrast
was repeated on seven RDKit axes computed directly from the frozen SMILES:
heavy-atom count (\texttt{GetNumHeavyAtoms}), molecular weight
(\texttt{Descriptors.MolWt}), Labute approximate surface area
(\texttt{CalcLabuteASA}), topological polar surface area
(\texttt{Descriptors.TPSA}), Crippen cLogP (\texttt{MolLogP}), rotatable-bond
count (\texttt{Descriptors.NumRotatableBonds}) and ring count
(\texttt{CalcNumRings}) \cite{rdkit2025}.  Docking-44 used all 12,651 primary
processed rows; DOCKSTRING used the same sorted 15,000-row complete-case sample
drawn without replacement with seed 71 as the molecular-weight controls.  Its
complete-row indices were SHA-256 fingerprinted.

For each descriptor separately, the empirical 25th and 75th percentiles were
calculated and boundary ties were retained with $z\leq q_{0.25}$ and
$z\geq q_{0.75}$.  Consequently, discrete descriptors could yield more than
25\% of the support in a tail; the actual low- and high-tail counts were
recorded and used for their matched controls.  Raw and row-centred target maps
were estimated on both extremes and compared by upper-triangle Spearman
correlation; sign-flip fraction and the PR of each extreme map were secondary
summaries.  All 21
pairwise Spearman correlations among the seven descriptors were also recorded
to expose their dependence.  The three labels ``size related'' (heavy-atom
count, molecular weight and Labute ASA) and ``other coarse'' were used only for
descriptive summaries.  For each distinct pair of low- and high-tail sizes, the
row-centred analysis also used 200 matched-size random-row-disjoint controls:
a seeded permutation supplied the first $N_{\mathrm{low}}$ rows to one map and
the next $N_{\mathrm{high}}$ rows to the other.  Size pairs were sorted
lexicographically and identical pairs were simulated once and reused.  The
generator seed was the dataset base plus $100{,}000$ times the zero-based size-
pair index; dataset bases were 20260821 for Docking-44 and 21260821 for
DOCKSTRING.  Control medians and empirical 2.5th and 97.5th percentiles were
descriptive finite-support comparators; no randomization probability was
assigned.  Because all seven axes were inspected post hoc and are correlated,
their contrasts are neither independent tests nor a causal decomposition of
the molecular-weight result.

\subsection{Pilot-map recovery and target-panel compression}

For each full processed source matrix, 100 simple-random pilot samples of 200
and 500 rows were drawn without replacement (seeds 20260810 for Docking-44 and
20261810 for DOCKSTRING).  The generator advanced continuously across the two
sample sizes.  Residual maps were calculated from row-centred scores.  For every
replicate, the row-disjoint complement map was computed exactly by subtracting
the pilot residual sum and cross-product from full-support sufficient
statistics; no approximation or resampling was used for the complement.

Pilot and reference maps were compared by upper-triangle Spearman correlation,
root-mean-square error and edge-sign agreement.  Strong-edge sign recovery used
all reference edges at or above the empirical 90th percentile of absolute
correlation.  Exact sets of the 10 and 25 most positive and most negative edges
were selected with stable mergesort tie handling and compared by overlap,
precision and Jaccard index.  Average-linkage clustering with optimal leaf
ordering used signed distance $1-r$ and was cut at $k=6,8,10$ for the cluster
sensitivity.

Panel compression treated either strong positive or strong negative linear
dependence as computational redundancy.  For a $k$-target panel $S$, the loss
on map $C$ was
\[
 L(S;C)=\frac{1}{P}\sum_{j=1}^{P}\min_{s\in S}\{1-|C_{js}|\}.
\]
The main decision used $k=8$; $k=6$ and $k=10$ were sensitivities.  Pilot maps
used deterministic PAM BUILD followed by SWAP, with target-name ordering to
break objective ties.  The fixed full-source reference panel was the exact
mixed-integer $k$-medoids solution with zero relative MIP gap and a negligible
index penalty used only to break exact ties.  Both panels were evaluated on the
row-disjoint complement.  The displayed difference is pilot-panel loss minus
the loss of this one fixed full-source panel; it is not regret against a panel
re-optimized on each complement.

For each dataset and each $k$, 10,000 uniformly sampled target subsets provided
the full-source random-panel comparator.  Their seeds were the dataset recovery
seed plus $10,000+k$ (20270818 and 20271818 for the main $k=8$ analyses).  To
limit repeated complement computation, the first 200 of those same fixed random
panels were evaluated on every complement.  Pilot summaries and their central
95\% ranges are empirical summaries of 100 samples from the observed library;
they neither establish a universal calibration size nor test transport to a
shifted library.

\subsection{Chemical-group-held-out pilot recovery}

To test whether within-source recovery depended on reusing chemical groups, a
separate sensitivity used five-fold \texttt{GroupKFold} with one harmonized
structure-derived grouping rule.  Each SMILES underwent RDKit cleanup,
largest-fragment selection, uncharging, isotope removal and removal of
stereochemistry.  Canonical non-isomeric Bemis--Murcko scaffold SMILES defined
cyclic groups.  Acyclic compounds were grouped by the first block of the
Standard InChIKey of the standardized parent, with canonical non-isomeric
SMILES as a declared fallback if InChI generation failed.  Docking-44 used all
12,651 rows and 5,112 groups; DOCKSTRING used the fixed 15,000-row complete-case
support drawn with seed 71 and 11,903 groups.  This Docking-44 grouping is
distinct from the frozen \texttt{Butina\_clusters} labels used in the
molecular-weight domain controls.  GroupKFold was deterministic and placed
every group wholly in one of five folds; the code asserted zero group overlap
between the four-fold calibration pool and the held-out fold.

Scores were clipped above at zero before splitting.  DOCKSTRING had no missing
score on this support.  For each Docking-44 split, target means were estimated
only from the four-fold calibration pool and applied to missing cells in both
that pool and its held-out fold.  Thus no held-out score informed imputation.
The held-out residual target map used every ligand in the fold.  From the
filled calibration pool, 25 simple-random samples without replacement were
drawn at each calibration size, 200 and 500; each row-centred map was compared
with the held-out-group map by upper-triangle Spearman correlation.

A matched random-row-holdout comparator used the same evaluation-fold size and
calibration size for every fold and repetition.  An independent permutation of
all rows supplied the random held-out rows first and the calibration pool
second; pool-only target means were again applied to both parts, and the first
$n$ pool rows formed the calibration sample.  Chemical groups were allowed to
overlap in this comparator.  Docking-44 used dataset base seed 20260805 and
DOCKSTRING used 20260806.  For one-based fold $f$ and calibration size $n$, the
chemical-group sample generator used
\texttt{dataset\_seed+10000*f+n}; the matched-random-row generator used that
seed plus 1,000,000.  Each generator advanced through the 25 replicates.

For every fold and design, achiral 2,048-bit radius-2 Morgan fingerprints were
computed on the same standardized parents.  The exact maximum Tanimoto
similarity of every held-out ligand to every ligand in its complete calibration
pool was recorded, requiring 25.6--36.0 million comparisons per fold.  The
group rule is deliberately stricter than a row split but is not a
fingerprint-cluster-disjoint split: a Tanimoto value of one denotes equality of
the declared finite bit vectors, not necessarily molecular identity, and
chemotypes with different group labels can remain similar.

The 125 values per dataset, design and calibration size were summarized by
their mean, sample standard deviation, extrema and empirical 2.5th and 97.5th
percentiles.  Paired chemical-group-minus-random-row differences used matching
dataset, fold, size and repetition indices and were summarized by their mean,
median and empirical 2.5th and 97.5th percentiles.  These ranges combine fixed-
fold composition with repeated calibration samples within one source
collection.  They are same-source holdout sensitivities, not confidence
intervals, evidence of transport to a shifted deployment library, an
equivalence test, or recovery of individual ligand scores.

\subsection{Chemical-group-held-out score-surface reconstruction}

A separate five-fold \texttt{GroupKFold} experiment tested whether a selected
target subset reconstructed held-out score surfaces. It used the same
standardized-parent chemical-group rule described above and asserted zero group
overlap. Only rows complete across all target columns entered the pilot and
evaluation sets. Within each fold, 500 pilot ligands were sampled from the four
training folds; the fifth fold remained unused for centring, scaling, target
selection, Ridge regularization or model fitting. The evaluation predictor
received only the observed raw Vina scores at the selected targets.

Panels of $k=4,6,8,10,12$ targets were selected by exact $k$-medoids on distance
$1-r^2$ from the pilot row-centred residual map. A pivoted-QR selector on the
pilot raw-score correlation representation and 50 fixed random panels per fold
and $k$ were controls. Multivariate Ridge regression predicted the omitted raw
targets; its penalty was selected from $\{0.01,0.1,1,10,100\}$ by generalized
cross-validation using pilot data only. Selected observed scores were copied
into the full predicted raw profile. That profile was then row-centred in raw
score units and standardized with pilot-only residual column scales to define
the derived residual prediction. Primary plots report variance-weighted
$R^2=1-\sum_j\mathrm{SSE}_j/\sum_j\mathrm{SST}_j$ on omitted targets; pooled
RMSE, target correlations, ligand-profile correlations, full-profile metrics,
a selected-target PC1 model and a seven-descriptor Ridge model were retained as
sensitivities.

At $k=8$, a direct-outcome sensitivity fitted the held-out residual columns
from the same selected raw target predictors, using the same pilot-only Ridge
procedure. Agreement with the derived-residual route tests whether low residual
reconstruction was mechanically introduced by predicting and then centring raw
scores. A designed-versus-random comparison was restricted to targets omitted
by both the designed panel and each of the first three fixed random panels; the
reported common-omitted contrasts are descriptive because only three random
panels were fixed for this robustness calculation.

\subsection{Experimental panels, exact ligand matching and map reliability}

The fixed target order was ABL1, AKT1, AKT2, CDK2, CSF1R, EGFR, FGFR1, IGF1R,
JAK2, KDR, KIT, LCK, MAP2K1, MAPK1, MAPK14, MAPKAPK2, MET, PLK1, PTK2, ROCK1
and SRC.  Source aliases were mapped explicitly rather than inferred from
strings.  DAVIS-Complete V3 was read from the exact CC0 TSV export
\path{data/frozen/davis_complete.tab.gz}; pivoting the selected proteins by
\texttt{drug\_name} produced a dense 72-by-21 pKd block \cite{davis2011,wu2025davis}.
PKIS2 was read from sheet \texttt{Table 4 - PKIS2 \%Inh} of the official S4
workbook \path{data/frozen/pkis2_s4.xlsx}.  Rows missing any of \texttt{Regno},
\texttt{Compound}, \texttt{Chemotype} or \texttt{Smiles} were excluded, leaving
645 rows; the 21 explicitly mapped percentage inhibition/displacement columns
were required to be numeric, dense and within $[0,100]$ \cite{drewry2017}.

Raw and row-centred target maps were computed independently within each assay
matrix. No ligand mean, target scale or other centring statistic was shared
between DAVIS, PKIS2 and HotSpot; cross-assay agreement therefore compares
separately estimated maps rather than a jointly normalized matrix.

The Anastassiadis HotSpot panel was read from Supplementary Table~3 of the
official publisher workbook \cite{anastassiadis2011}. The fetcher records the
source URL and accepts the file only at its frozen SHA-256; the workbook is not
redistributed. Percentage remaining activity was converted to
$100-\text{percentage remaining}$ without clipping. CDK2 was the per-compound
arithmetic mean of the cyclin-A and cyclin-E rows. Of 178 compounds, SB~202474
was excluded for a missing AKT1 value and VEGF Receptor~2 Kinase Inhibitor~II
for a missing MAPK14 value, leaving a dense \HotspotLigands{}-by-21 block.
Target aliases, worksheet cells, exclusions and alternative CDK2 contexts are
serialized in the result bundle.

Profiles with range no greater than $10^{-12}$ across targets were called
constant; the informative cross-panel maps therefore used 67 DAVIS and 644
PKIS2 profiles, whereas the primary split-half summaries used every released
row (72 and 645, respectively).  The two-way-centred map was primary.  Raw maps
were repeated for split-half map reliability and for same-support point and
bootstrap comparisons; edge-stability and paired same-support split analyses
used the primary transform.

Molecular identities for DAVIS and PKIS2 were recomputed as Standard InChIKeys
from the source SMILES with pinned RDKit.  The complete DOCKSTRING support was
aligned to the frozen raw-identity contract described above.  Exact matching
required equality of the full Standard InChIKey.  If an identity occurred more
than once on either side, each ligand--target cell was collapsed by the median
within that side before the sorted identity intersection was aligned.  Vina
scores in the matched block were clipped above at zero.  This yielded 59 DAVIS
exact identities, of which three had a constant experimental profile at
pKd=5, giving the primary 56-ligand informative block; all 154 PKIS2 exact
identities were informative.

The broad docking reference retained complete DOCKSTRING rows after excluding
every row whose raw 14-character InChIKey connectivity block occurred anywhere
in either full DAVIS or full PKIS2 panel.  This removed 254 rows and retained
259,806.  Its positive scores were clipped at zero.  The same 21 targets and
two-way-centred map definition were then used for the broad-versus-experimental
comparison.  The exact-support and broad-reference comparisons therefore
differ only in ligand-support contract, not target set or transformation.

Chemical clusters for all experimental split and bootstrap analyses were
generated from binary RDKit Morgan fingerprints using
\texttt{GetMorganGenerator(radius=2, fpSize=2048)} and
\texttt{GetFingerprint} \cite{morgan1965,rogers2010,rdkit2025}.  Pairwise
distance was one minus Tanimoto similarity.  Butina clustering used a distance
cutoff of 0.35 (Tanimoto similarity at least 0.65),
\texttt{isDistData=True} and \texttt{reordering=True} \cite{butina1999}.  For
completeness, the public loader's Murcko labels used canonical cyclic scaffolds
and grouped an acyclic molecule by its Standard-InChI connectivity block
(\texttt{ACYCLIC:<connectivity>}); those Murcko labels were not the grouping
unit in the reported experimental splits or bootstraps.  This differs from the
source-row-singleton convention used in the DOCKSTRING domain control.

Disjoint-half repeatability used 2,000 repetitions.  A ligand split randomly
permuted rows and assigned $\lfloor N/2\rfloor$ rows to each half, leaving one
row unused when $N$ was odd.  A cluster split randomly ordered complete Butina
clusters and assigned each successive cluster to the currently smaller half;
no cluster crossed halves.  Each half was mapped independently.  One complete
target-label relabeling per split supplied local sign and top-decile null
comparators.  Edge perturbation used 2,000 ordinary ligand bootstraps and 2,000
Butina-cluster bootstraps.  A cluster bootstrap sampled the number of distinct
clusters with replacement and included every member of each sampled cluster,
so its ligand count could vary.

Same-support uncertainty used paired $n$-out-of-$n$ resampling: the same ligand
or Butina-cluster draw was applied to the matched docking and experimental
matrices.  The ordinary empirical-correlation result was primary; Oracle
Approximating Shrinkage covariance converted to correlation was a finite-sample
estimator sensitivity \cite{chen2010}.  A separate paired split-half analysis
compared within-modality and cross-half maps directly.  No Spearman--Brown
extrapolation, attenuation correction or classical noise ceiling was applied,
because the nonlinear map statistic and the two assay modalities cannot be
assumed to measure one interchangeable latent quantity.

The experimental DAVIS--PKIS2 comparison used the informative 67- and
644-profile maps.  For each of 9,999 QAP permutations, one complete permutation
of the PKIS2 target labels was applied jointly to its rows and columns while the
DAVIS map remained fixed \cite{hubert1976}.  The two-sided probability was
$(1+\#\{|\rho_b|\geq|\rho_{\mathrm{obs}}|\})/(9,999+1)$.  With 210 target-pair
edges, the top positive decile contained $\lceil0.10\times210\rceil=21$ edges,
selected with deterministic stable tie handling.  QAP respects the shared-node
dependence of map edges, but remains conditional on this fixed target panel.

HotSpot repeatability used 2,000 disjoint random ligand-row halves because the
publisher workbook contains names and CAS numbers but not structures. A
separate PubChem identity audit resolved 155 of 176 records; unresolved records
were retained as singletons. The resolved standardized structures defined
radius-2, 2,048-bit Morgan fingerprints and Butina clusters at Tanimoto 0.5 for
a secondary positive-multiplier sensitivity, but not for the primary split-half
summary. Pairwise cross-assay concordance compared corresponding upper-triangle
edges of the complete DAVIS, PKIS2 and HotSpot raw or row-centred maps.

All experimental resampling used base seed 20260810, 2,000 split repetitions
and 2,000 bootstraps.  DAVIS used panel seed 20260810 and PKIS2 used 21260810
(panel stride 1,000,000).  Split seeds were
\texttt{panel\_seed+10000+counter} in the serialized loop order over full then
informative support, two-way then raw transform, and ligand then Butina unit.
For the edge bootstraps, \texttt{counter} was 4 after the full-support splits
and 8 after the informative-support splits, and \texttt{unit} was 0 for ligand
and 1 for Butina resampling; the seed was
\texttt{panel\_seed+100000+10*counter+unit}.  Paired same-support bootstrap seeds
were
\[
\texttt{panel\_seed}+200000+10000\,\texttt{support}
 +100\,\texttt{transform}+\texttt{unit},
\]
where \texttt{support}=0 for the primary informative exact matches and 1 for all
exact matches, \texttt{transform}=0 for two-way and 1 for raw, and the same unit
indices.  Paired same-support split seeds were
\[
\texttt{panel\_seed}+300000+10000\,\texttt{support}+\texttt{split},
\]
with \texttt{split}=0 for ligand and 1 for Butina.  Cross-panel QAP seeds were
$20260810+5{,}000{,}000=25260810$ for two-way-centred and 25260811 for raw
maps.  Percentile bootstrap and repeated-split ranges are
sampling/perturbation ranges for nonlinear statistics on observed supports, not
target-population confidence intervals.

\subsection{Assay-value-blind fixed-target panel transfer}

Panel transfer used the same ordered 21 targets and a Vina calibration matrix
that excluded the union of all DAVIS and PKIS2 connectivity blocks and every
conservative HotSpot identity candidate. Resolved HotSpot structures contributed
their audited RDKit connectivity; each unresolved record contributed all
plausible PubChem candidate connectivities. This removed 419 rows from the
260,060-row complete DOCKSTRING support and retained 259,641 rows. Positive
Vina scores were clipped at zero, and row centring used only the fixed 21 target
columns. Identity resolution and exclusion were performed without assay values.

For $k=4,6,8,10,12$, all $\binom{21}{k}$ panels were enumerated. Separate raw
and row-centred Vina maps selected every numerically optimal panel under each of
three distances: signed $1-r$ (primary), $1-|r|$, and $1-r^2$. Each selected
panel was then evaluated on the corresponding row-centred DAVIS, PKIS2 or
HotSpot experimental map using the same coverage loss, the mean nearest-target
distance from all 21 targets to the selected set. No experimental value entered
selection. Ties were reported lexicographically, averaged over all optima, and
conservatively as the best raw-panel loss minus the worst residual-panel loss.
Positive raw-minus-residual contrast favours the residual selector.

The primary joint summary was the mean across the three assay panels of the
normalized trapezoidal area under the contrast $c(k)$ over $k$. For the ordered
grid $k_1<\cdots<k_M$, this was
\begin{equation*}
\mathrm{AUC}_{\mathrm{norm}}
=\frac{1}{k_M-k_1}\sum_{m=1}^{M-1}
\frac{(k_{m+1}-k_m)\{c(k_m)+c(k_{m+1})\}}{2},
\end{equation*}
so normalization is by the observed $k$ range only and leaves the signed
$1-r$ loss scale unchanged. One common random
permutation of all 21 target labels was applied to every assay panel, $k$,
objective and tie statistic in each of 20,000 repetitions. One-sided
probabilities used the add-one rule. A maximum-statistic adjustment over the
three distance objectives was reported separately for each tie statistic; the
conservative signed result is primary. A sequence baseline applied the same
exact panel enumeration to the fixed 21-target global sequence-identity map.
Re-enumerating a selector after relabelling its complete Vina map is exactly
equivariant to applying that relabelling to every member of the already
enumerated numerical-optimum set. The implementation uses the latter form and
then evaluates those permuted optima against each fixed experimental distance
map; selection is therefore inside the target-label null rather than held fixed
in biological label space.

For the HotSpot chemical-support sensitivity, every member of a resolved
Morgan/Butina cluster received one shared independent $\mathrm{Exp}(1)$ weight;
the 21 unresolved compounds were singleton clusters. Two thousand weighted
maps were evaluated with the already selected Vina panels. This multiplier is a
descriptive perturbation of chemical support; target-label permutation remains
the inferential unit for panel transfer.

Target influence was assessed by omitting each kinase in turn, recomputing all
raw and row-centred maps, and repeating complete panel enumeration on the
remaining 20 targets for the same $k$ grid and primary signed objective. Thus
the selectors themselves, rather than only their evaluation edges, changed in
each leave-one-target-out analysis. Cross-assay raw and residual map agreement
was recomputed on the same reduced panels. The resulting ranges are fixed-panel
influence sensitivities, not intervals for a target superpopulation.

\subsection{Docking-44 missing-data analyses}

The primary target-mean imputation was compared with target-median imputation
on the same 12,651 rows and with a complete-case analysis on the 9,297 rows
having all 44 scores.  Restricting the primary imputed matrix to these rows is
exactly identical to their observed score matrix, making the complete-case
contrast a support contrast rather than an imputation contrast.

An observed-cell additive least-squares sensitivity estimated grand, ligand and
target effects by alternating exact unweighted updates over available cells.
The convergence tolerance was $10^{-12}$ with at most 10,000 iterations; the
released fit converged in seven.  Residuals were calculated at observed cells,
and every missing residual was set to zero, its fitted additive expectation.
This completes a residual surface for spectral calculation; it does not impute
or recover the missing raw docking scores.

To separate the unusual complete-row chemistry from its smaller sample size,
1,000 simple-random supports of 9,297 rows were drawn without replacement from
the full primary processed surface (seed 20260811).  Their residual PR
distribution is a descriptive equal-$N$ comparator, not a randomization test of
missingness.  Molecular weight, rotatable-bond count and cLogP were read from
the frozen \texttt{MW, g/mol}, \texttt{rotatable\_bonds} and \texttt{logP}
columns; heavy-atom count was recomputed with RDKit from cleaned SMILES,
falling back to canonical SMILES.  These quantities were summarized for
complete and incomplete rows, with excluded-minus-complete raw shifts,
standardized mean differences, median/IQR shifts and Cliff's delta.  No missing-
at-random mechanism was assumed and no hypothesis test assigned causality to
these associations.

\subsection{Software and computational reproducibility}

The public-core environment used Python~3.12.11, NumPy~1.26.4,
pandas~2.3.2, SciPy~1.16.2, scikit-learn~1.7.2, RDKit~2025.03.6,
Matplotlib~3.10.6 and openpyxl~3.1.5; pytest~9.1.1 ran the focused tests
and xlrd~2.0.2 read the publisher XLS
\cite{numpy2020,scipy2020,rdkit2025,matplotlib2007}.  The exact environment is
pinned in \path{requirements.lock} and \path{Dockerfile}.

The source-level public producers are:
\begin{itemize}
\item \path{analysis/cross_panel_overlap_audit.py};
\item \path{analysis/public_residual_null_audit.py};
\item \path{analysis/public_docking44_missing_data_sensitivity.py};
\item \path{analysis/public_chemical_domain_controls.py};
\item \path{analysis/public_descriptor_domain_specificity.py};
\item \path{analysis/revision_map_decision_sensitivities.py};
\item \path{analysis/public_scaffold_holdout_recovery.py};
\item \path{analysis/public_panel_selector_controls.py};
\item \path{analysis/public_panel_selector_robustness.py};
\item \path{analysis/experimental_map_reliability.py};
\item \path{analysis/fetch_anastassiadis_supplement.py};
\item \path{analysis/public_anastassiadis_identity_audit.py}; and
\item \path{analysis/public_anastassiadis_panel_validation.py}.
\end{itemize}
Running each with its recorded defaults recreates its named result directory
from the redistributed row-level sources, the checksum-fetched HotSpot workbook
and the checksum-listed public identity and target annotations. Corresponding
\texttt{test\_*.py} files lock dimensions, preprocessing, support fingerprints,
singleton handling, identity matching, seed contracts and numerical invariants.

Finally, \path{analysis/build_public_core_evidence.py} reopens the allowlisted
public sources and producer outputs, recomputes the primary spectra directly,
checks cross-artifact consistency, and writes the schema-4.0.0 ledger
\path{results/public_core_evidence.json}.  It also writes
\path{results/public_core_source_manifest.csv} and its SHA-256 digest; every
opened input and output is recorded with path, role, byte size, license and
SHA-256.  The independent \path{analysis/target_map_audit.py} command-line tool
fails closed on missing or nonnumeric target cells unless target-mean or
target-median imputation is explicitly requested, requires an absent or empty
output directory, and emits the preprocessing contract, raw and two-way maps,
spectra, signed edges and a per-run checksum manifest.  Its optional domain,
row-disjoint pilot and deterministic $1-|r|$ medoid analyses implement the
study's reusable audit; they do not validate biological mechanism, affinity
accuracy or ligand-wise target retrieval.
